% Original: PDF strana 7, gornji slajd.
\slajdid{02-s013}
\begin{frame}[label=02-s013]{Zbir proizvoda iz redova sa \(F=1\)}
% element: 02-s013:e04
Na osnovu osobina logičkih funkcija, iz funkcionalne tabele možemo napisati izraz za logičku funkciju.\par\vspace{2mm}
\begin{columns}[T,onlytextwidth]
\begin{column}{.62\textwidth}
% element: 02-s013:e01

\[\begin{aligned}
(C,B,A)=(0,1,1)&\Rightarrow\bar CBA\\[2mm]
(1,0,1)&\Rightarrow C\bar BA\\[2mm]
(1,1,0)&\Rightarrow CB\bar A
\end{aligned}\]
\par\vspace{2mm}
% element: 02-s013:e02
\Oznaka{Veze između uslova}\vspace{2mm}
Unutar jednog reda literali su povezani sa \alert{I}; alternativni redovi sa \alert{ILI}. Nula u redu daje komplementni literal.
\par\vspace{3mm}
% element: 02-s013:e03
\Rezultat{\(\textcolor{DEcrvena}{F=\bar CBA+C\bar BA+CB\bar A}\)}
\end{column}
\begin{column}{.32\textwidth}\centering
\Oznaka{Funkcionalna tabela}\vspace{3mm}
{\fontsize{9}{11}\selectfont \begingroup
\setlength{\tabcolsep}{3pt}
\setlength{\aboverulesep}{1pt}\setlength{\belowrulesep}{1pt}
\renewcommand{\arraystretch}{1.05}
\par\noindent
\begin{tabularx}{\linewidth}{*{4}{>{\centering\arraybackslash}X}}
\toprule C&B&A&F\\\midrule
0&0&0&0\\\cmidrule[.2pt](lr){1-4}
0&0&1&0\\\cmidrule[.2pt](lr){1-4}
0&1&0&0\\\cmidrule[.2pt](lr){1-4}
\textcolor{DEcrvena}{0}&\textcolor{DEcrvena}{1}&\textcolor{DEcrvena}{1}&\textcolor{DEcrvena}{1}\\\cmidrule[.2pt](lr){1-4}
1&0&0&0\\\cmidrule[.2pt](lr){1-4}
\textcolor{DEcrvena}{1}&\textcolor{DEcrvena}{0}&\textcolor{DEcrvena}{1}&\textcolor{DEcrvena}{1}\\\cmidrule[.2pt](lr){1-4}
\textcolor{DEcrvena}{1}&\textcolor{DEcrvena}{1}&\textcolor{DEcrvena}{0}&\textcolor{DEcrvena}{1}\\\cmidrule[.2pt](lr){1-4}
1&1&1&0\\\bottomrule
\end{tabularx}\par
\endgroup
}
\end{column}
\end{columns}
\end{frame}
